% !TeX program = xelatex
% !TeX encoding = UTF-8
\documentclass[UTF8,zihao=-4,fontset=none,a4paper]{ctexart}
% 编译建议：XeLaTeX 或 Tectonic。请保留同目录下的 fonts 文件夹。
% 本地嵌入 Noto Serif SC，避免依赖阅读器额外提供 Adobe-GB1 中文映射。
\setCJKmainfont[Path=fonts/,BoldFont=NotoSerifSC-Bold.otf]{NotoSerifSC-Regular.otf}
\setCJKsansfont[Path=fonts/,BoldFont=NotoSerifSC-Bold.otf]{NotoSerifSC-Regular.otf}
\setCJKmonofont[Path=fonts/,BoldFont=NotoSerifSC-Bold.otf]{NotoSerifSC-Regular.otf}
% 分工为开题阶段的建议安排，后续按小组实际工作更新。
\usepackage[left=2.1cm,right=2.1cm,top=1.9cm,bottom=1.9cm]{geometry}
\usepackage{amsmath,amssymb,bm}
\usepackage{booktabs,tabularx,array}
\usepackage{enumitem}
\usepackage{fancyhdr}
\usepackage[hidelinks]{hyperref}
\usepackage{microtype}
\hypersetup{
  pdftitle={胚胎发育时间外推中细胞组成与表达变化的统计分析},
  pdfauthor={统计学习课程项目小组},
  pdfsubject={Virtual Embryo Challenge T1 开题报告}
}
\linespread{1.12}
\setlength{\parindent}{2em}
\setlength{\parskip}{2pt}
\setlength{\headheight}{15pt}
\setlength{\emergencystretch}{1em}
\setlist{nosep,leftmargin=2em}
\ctexset{section={format=\large\bfseries,beforeskip=9pt,afterskip=4pt},
  subsection={format=\normalsize\bfseries,beforeskip=5pt,afterskip=2pt}}
\pagestyle{fancy}
\fancyhf{}
\fancyhead[L]{\small 2026 秋季《统计学习》课程大作业}
\fancyhead[R]{\small 开题报告}
\fancyfoot[C]{\small \thepage}
\renewcommand{\headrulewidth}{0pt}
\renewcommand{\footrulewidth}{0pt}
\newcommand{\smalltable}{\small\renewcommand{\arraystretch}{1.18}}
\newcolumntype{Y}{>{\raggedright\arraybackslash}X}
\begin{document}
\begin{center}
{\zihao{3}\bfseries 胚胎发育时间外推中细胞组成与表达变化的统计分析\par}
\vspace{5pt}
{\normalsize NeurIPS 2026 Virtual Embryo Challenge T1 Temporal\par}
\vspace{4pt}
{\small 小组成员：高景行\quad 王恒毅\quad 李凯瑞\quad 崔帅达\par}
\vspace{2pt}
{\small 日期：2026 年 10 月 18 日}
\end{center}
\vspace{-5pt}

\section{问题背景与研究范围}
胚胎发育伴随细胞分化、群体比例改变和基因表达重组。单细胞 RNA 测序能够观测不同发育阶段的表达状态，但分时采样缺少逐细胞的前后对应关系。已有研究使用最优传输分析时间快照并推断分化过程\cite{schiebinger}，也提示这类问题需要在模型假设下研究群体演化。

本组拟聚焦 Virtual Embryo Challenge 的 T1：根据较早阶段数据，预测未来阶段的一组细胞表达向量\cite{task}。预测目标是细胞群体的分布，包括状态组成及表达结构；本项目不将其解释为个体细胞的真实发育轨迹。研究重点是建立可解释的统计模型，比较其外推能力并分析失败原因。

\section{具体研究问题}
核心问题是：\textbf{在观测时间点有限的条件下，低维且受约束的外推能否改善未来细胞群体预测？} 拟进一步研究：
\begin{enumerate}[label=（\arabic*）]
  \item 降维与收缩能否减轻高维表达变化外推时的噪声放大？
  \item 调整细胞群体比例、调整群体内部表达分别带来哪些预测收益？
  \item 方法的表现是否随细胞群体、降维维数和外推强度发生变化？
\end{enumerate}
以上均为待检验问题。主要工作定位为统一条件下的模型比较、消融与误差分析，结果将同时报告改善、退化及不稳定的情形。

\section{数据来源与使用方案}
使用官方发布的 T1 单细胞 RNA 数据。数据以 AnnData 文件组织，表达矩阵已作 log1p 归一化，包含 32,285 个基因及训练阶段的细胞类型注释\cite{data}。下载后将检查文件版本、矩阵含义、缺失值、基因顺序和注释一致性，避免重复归一化；特征处理和 PCA 仅在相应训练阶段上拟合。

\begin{center}
\smalltable
\begin{tabularx}{\linewidth}{@{}p{2.0cm}p{2.8cm}Y@{}}
\toprule
角色 & 发育阶段 & 研究用途\\
\midrule
初期训练 & E8.5、E9.5 & 数据探索、表示学习与变化规则估计\\
开发验证 & E10.5 & 按发布安排评价预测及选择模型\\
隐藏测试 & E12.5 & 最终官方评价，真值不公开\\
\bottomrule
\end{tabularx}
\end{center}

T1 来自以心脏为中心的组织取样，阶段间取样范围与标签命名可能不同\cite{data}，因此须先检查可比性，再解释群体变化。本组已完成官方平台团队注册，并确认可下载 \texttt{E8.5\_RNA.h5ad}、\texttt{E9.5\_RNA.h5ad}；数据内容将在获取后按上述流程校验。

\newpage
\section{初步建模思路}
\textbf{参考基线。} 复现官方 \texttt{copy\_last} 和 \texttt{pseudobulk\_shift}\cite{task}，分别沿用最近阶段分布，以及按已观察的平均表达变化整体平移。另设\textbf{全基因空间收缩外推}：在下式中将 $V_dV_d^\top$ 替换为单位矩阵，保持时间倍率、收缩强度、截断和抽样一致，以区分降维与收缩的作用。

\textbf{低维收缩外推。} 设训练阶段为 $t_a<t_b$，在其数据上拟合 PCA，得到基矩阵 $V_d\in\mathbb{R}^{p\times d}$。记阶段均值为 $\bar{\bm x}_a,\bar{\bm x}_b$，时间倍率 $h=(t_\star-t_b)/(t_b-t_a)$。对最近阶段抽取的细胞，预测为
\begin{equation}
\widehat{\bm x}_i(t_\star)=
\left[
\bm x_i(t_b)+
\alpha h\,
V_dV_d^\top(\bar{\bm x}_b-\bar{\bm x}_a)
\right]_+ ,
\label{eq:model}
\end{equation}
其中 $\alpha\in[0,1]$ 为收缩强度，$[\cdot]_+$ 为逐元素非负截断；$\alpha=0$ 对应最近阶段基线。模型保留抽样细胞的高维残差，并记录截断比例。PCA 可能遗漏时间变化方向，须与全空间收缩模型在相同 $\alpha$ 下比较。

\textbf{分层扩展与消融。} 先依据训练资料确定可靠的共同标签或较粗分类；仅对两阶段各不少于 30 个细胞的匹配群体，在同一 PCA 基上估计群体内均值差。令 $\bm\pi_a,\bm\pi_b$ 为匹配集合内部的条件比例，按
$\widehat{\bm\pi}=\operatorname{Normalize}([\bm\pi_b+\beta h(\bm\pi_b-\bm\pi_a)]_+)$
更新，$\beta\in[0,1]$，Normalize 表示除以分量之和。匹配集合的总质量及其他群体比例保持为最近阶段值，再按预测比例重抽样；未匹配或样本不足群体的表达也沿用最近阶段。若无合格群体则回退至整体模型。比较“均不调整”“仅调整表达”“仅调整比例”“同时调整”四种设置。

初设 $d\in\{10,30,50\}$，$\alpha,\beta\in\{0,1/2,1\}$，保留全部比较记录。共同均值平移不能主动学习群体内协方差演化；消融仅比较建模选择的经验作用，不作因果分解。

\section{评价方案与泄漏防范}
采用按时间划分的\textbf{开发验证}。官方计划于 10 月 20 日公开 E10.5 真值\cite{timeline}；公开后仅用 E8.5、E9.5 拟合预测模型，在 E10.5 上评价并作有限模型选择，保存配置、预测文件及 Git 版本。E10.5 参与调参，相关结果均标记为开发验证结果，不作为独立测试。同阶段随机拆分仅用于流程检查。

冻结方法、超参数和标签映射规则后，按官方许可用 E9.5、E10.5 重新拟合 PCA 与变化率，取 $h=2$ 预测 E12.5。一步开发验证不能保证两天外推有效。正式提交前先做格式检查，并遵守测试阶段全阶段至多两次官方评分的现行限制，不根据评分反解目标信息\cite{rules}。

主评价采用固定版本的官方工具，报告 DES（差异表达恢复，25\%）、DCS（变化方向，25\%）、MMD（细胞状态分布，30\%）及 CSS（基因共同变化，20\%）的原始值与官方总分；总分按官方 skill 校准后加权，不直接加权原始指标\cite{evaluation}。提交完整表达并按目标 board 的基因清单排序；细胞类型由官方冻结分类器判定。

补充分析包括群体误差、代表性基因图和参数敏感性。各方法保持一致的预测细胞抽样预算，并用固定的若干随机种子重复抽样，描述算法和抽样波动；这些重复及大量细胞均不视为独立胚胎重复。

\newpage
\section{可能困难与应对}
\begin{itemize}
  \item \textbf{时间点少与新状态出现。} 采用低自由度外推并分析失效群体；对未观察到的细胞状态保留预测局限，不声称恢复完整发育动力学。
  \item \textbf{取样与注释差异。} 根据训练资料检查标签对应，必要时在可靠的较粗层级比较；不将所有阶段差异归因于分化。
  \item \textbf{高维计算与输出约束。} 使用稀疏读取、适量抽样及低维计算；输出时核对非负性、有限值和官方基因顺序。
  \item \textbf{外部信息重叠。} 初期仅使用官方允许的训练文件。后续若引入外部数据或预训练模型，将核查训练阶段、基因型与来源并完整披露，遵守 held-out data 规定\cite{rules}。
\end{itemize}

\section{进度安排与初步分工}
已完成团队注册与训练数据下载和开题。10 月 20 日起开展开发验证，至 11 月 19 日完成数据探索、基线与初步模型；11 月下旬完成消融和误差分析，12 月 2 日前完成提交；随后准备汇报。以下为初步分工建议，后续按实际协作调整。

\begin{center}
\smalltable
\begin{tabularx}{\linewidth}{@{}p{1.8cm}p{2.5cm}Y@{}}
\toprule
成员 & 学号 & 初步职责\\
\midrule
高景行 & 2400010759 & 官网上数据获取、数据预处理与基本数据分析\\
王恒毅 & 2400010616 & 官网注册团队、参考Baseline、预测文件生成与官方提交流程\\
李凯瑞 & 2400010604 & 做低维外推、分层扩展实验与消融实验\\
崔帅达 & 2400010883 & 评价实现代码与误差解释\\
\bottomrule
\end{tabularx}
\end{center}
全组共同核查方法、准备报告与答辩，并依据实际工作更新贡献。预期交付可复现的统一比较、消融与失败案例分析；即使改进未超过基线，也将据实解释适用条件和局限。

\section{复现与 AI 辅助说明}
代码拟包含 README、数据获取说明、分析与评价脚本、环境版本和随机种子，并记录官方工具版本、配置与预测文件。项目的想法有与ChatGPT的交流和讨论；文档基本由组员完成，部分用词用ChatGPT润色。拟按 Human Team 工作方式开展实验，组员讨论并选择方法、检查结果及决定调整\cite{rules}。

\begingroup
\small
\setlength{\parskip}{0pt}
\begin{thebibliography}{9}
\setlength{\itemsep}{1pt}
\bibitem{task}
Virtual Embryo Challenge. \href{https://virtualembryo.ai/challenge/tasks/temporal}{Task 1 Temporal}.
\bibitem{data}
Virtual Embryo Challenge. \href{https://virtualembryo.ai/challenge/data}{Data}.
\bibitem{timeline}
Virtual Embryo Challenge. \href{https://virtualembryo.ai/challenge/timeline}{Timeline}.
\bibitem{rules}
Virtual Embryo Challenge. \href{https://virtualembryo.ai/challenge/rules}{Official Rules}.
\bibitem{evaluation}
Virtual Embryo Challenge. \href{https://virtualembryo.ai/challenge/evaluation?section=scoring\&task=1}{Task 1 Evaluation}.
\bibitem{schiebinger}
Schiebinger G, et al. Optimal-Transport Analysis of Single-Cell Gene Expression Identifies Developmental Trajectories in Reprogramming.
\textit{Cell}, 2019, 176: 928--943.e22.
\href{https://doi.org/10.1016/j.cell.2019.01.006}{doi:10.1016/j.cell.2019.01.006}.
\end{thebibliography}
\endgroup
\end{document}
